\section{Conservation determines every weight on a certified single orbit}
\label{sec:single-orbit-weights}

Normal-component queries attach additive integer vectors to interval points:
Euler characteristic and boundary intersection numbers for disc detection,
or normal-coordinate counters for component extraction. The maintained AHT
producer first computes an unweighted orbit certificate, then transports
these weights along its reductions. Once that certificate proves that there
is exactly one orbit, the transport step has a simpler exact replacement.

\subsection{A complete invariant under a checked connectivity hypothesis}

Let $X=\{0,\ldots,N-1\}$ and let $w:X\to\mathbb Z^d$ be the additive
weight function. If the interval-pairing relation has exactly one orbit,
its histogram is forced to be
\[
 \left\{\left(\sum_{x\in X}w(x),\ 1\right)\right\}.
\]
For a canonical run partition $(\ell_i,r_i,v_i)_{i=1}^R$, using half-open
intervals, the weight is computed without expanding points:
\[
 W=\sum_{i=1}^R (r_i-\ell_i)v_i.
\]
Signed coordinates, zero runs and overlapping source intervals are allowed;
the existing independent canonicalization first recovers their summed
piecewise-constant function. This is conservation plus \emph{proved}
connectedness, not a connectedness conjecture about primitive normal surfaces.

The producer takes this path only after its complete unweighted AHT run
returns one orbit. Public reuse of a supplied trace verifies it before
entering the same result routine. The independent weighted checker first
reconstructs the source weights, checks their binding to the certificate,
and fully replays the unweighted orbit proof. Only then may it replace
weight transport by the mass formula. A forged orbit count of one never
bypasses the local-rule checker.

The existing \code{weighted-interval-orbits-v1} format is sufficient. Its
canonical weights, unweighted proof and histogram are unchanged, so old and
new complete certificates agree exactly. This is shorter computation, not
a claim of fewer serialized certificate bytes. The producer uses direct
coordinate accumulation; the independent checker uses its own tuple-vector
arithmetic and reconstructed weight function. Neither implementation calls
the other's transport code to justify the shortcut.

For $R$ runs and dimension $d$, the extra mass calculation uses $O(Rd)$
integer products and additions and $O(d)$ accumulator space. If endpoint
lengths have at most $B$ bits and entries at most $C$ bits, intermediate
signed sums have $O(B+C+\log(R+1))$ bits. A conservative bit bound is
\[
 O\!\left(Rd\,\mathsf M(B+C+\log(R+1))\right),
\]
where $\mathsf M$ bounds integer multiplication and dominates addition.
This excludes source canonicalization and the mandatory unweighted proof
search or replay. It removes weighted work depending on all intermediate
transport partitions, but does not make the full orbit query linear.

\subsection{Why several orbits still require transport}

Conservation alone is insufficient when there is more than one orbit.
For example, pair $0$ with $2$ and $1$ with $3$, and assign weights
$1,2,3,4$ respectively. The two orbit weights are $4$ and $6$. A forged
histogram with two weights equal to $5$ conserves total mass and has the
right orbit count, but is false. The maintained multi-orbit path therefore
keeps the complete independent weighted transport and rejects that forgery.
An empty universe also retains its existing empty histogram; a one-point
universe uses the single-orbit identity.

For a connected normal surface the three-weight disc query now obtains
Euler characteristic and both boundary parities directly from total mass,
after authenticating connectedness by its interval proof. The existing
surface classifier and source-to-exterior verifier still supply the
topological and provenance obligations. No new genus assumption, negative
knot verdict or source-certificate schema is introduced. The same shortcut
works for connected full-coordinate queries and for signed abstract weights.

\subsection{Compatibility and source audits}

Four focused tests pass in 0.004 seconds. They exercise signed overlapping
weights on a $2^{20000}$-point connected system with both weight-transport
routines disabled, including JSON hexadecimal round trips and supplied-trace
reuse. They reject a false multi-orbit histogram with correct total mass,
a forged single-orbit count, modified unweighted operations, incorrect source
weights and Boolean counts. Empty and singleton universes retain their
expected behavior. Cancellation still propagates through verification.
The sixteen existing weighted tests also pass in 1.251 seconds.
The first full run exposes one obsolete cancellation threshold: the
six-tetrahedron fixture's verifier now finishes in 348 callbacks, before
the test's hard-coded 600th callback. The test now measures each operation
and requires cancellation at its entry, middle and final guards. All
seventeen component tests pass in 0.917 seconds; the original failure log
is retained. No artificial work is added to preserve a historical callback
count.
The final full suite passes all 1,175 tests in 242.337 seconds. The final
source audit takes 32.626 seconds while that suite also runs; this elapsed
time is not a performance comparison.

The pinned audit compares 300 random weighted systems with the published
producer and checker. All entire certificates agree. The old checker accepts
every new certificate and the new checker accepts every old certificate.
Of these systems, 226 have one orbit; only those use the new branch. Total
producer weight-replay events fall from 17,415 to 1,706. The latter events
remain necessary for the other systems; this count does not include the
mandatory unweighted proof operations.

All 84 source queries are repeated at the existing shared two-million-guard
allowance. The same seventeen positive proofs replay independently with
their previous hashes. Sixty-four restricted misses finish and three cases
remain capped, compared with sixty-three misses and four caps before.
The newly finished query is \code{survivor-02}, at 1,998,613 guards; it still finds no disc in
either candidate, so there is no new unknot verdict. Gordian and the
32-crossing stabilized circle remain among the capped cases. The 16-crossing
circle's shared work falls from 988,196 to 782,613 guards while preserving
its certificate. Budget counters measure cooperative work, not bit time.

\subsection{Complete weighted queries and source recognition}

The final measurements run serially after the successful full suite, final
audit and runtime publication. Each experiment uses five shuffled paired
rounds, old and new A/A controls and one excluded warmup per arm and input.
All 180 measured calls and 36 warmups finish. Every repeated old and new
certificate hash agrees within its input. Tables report separate-arm median
times and medians of within-round ratios.

The weighted endpoint constructs its input, computes the entire histogram
and certificate, independently verifies that certificate and computes its
hash. Connected cases use the earlier five-cycle interval family with 32
signed weight runs and either three or 128 coordinates. One case uses
4096-bit endpoints. A reflection with 128 two-point orbits is the
disconnected control. Both producer and verifier switch together between
the baseline and new implementations; the unweighted orbit engine is fixed.

\begin{center}
\small
\begin{tabular}{lrrrrr}
\toprule
Input & Old (ms) & New (ms) & Old/new & Old A/A & New A/A \\
\midrule
Connected, 3 weights & 4.840 & 3.930 & 1.188 & 0.947 & 1.055 \\
Connected, 128 weights & 13.064 & 7.648 & 1.729 & 1.003 & 0.977 \\
128 weights, 4096 bits & 31.456 & 20.706 & 1.542 & 1.004 & 1.023 \\
Disconnected control & 0.824 & 0.822 & 1.005 & 0.979 & 1.010 \\
\bottomrule
\end{tabular}

\end{center}

All eighty measured queries and sixteen warmups finish in 1.052 seconds of
driver elapsed time. Connected-case paired gains range from $1.188$ to
$1.729$; the 128-coordinate small-endpoint case falls from 13.064 to
7.648 milliseconds. The disconnected ratio is $1.005$, within the observed
controls. Old A/A ratios range from $0.947$ to $1.004$ and new A/A from
$0.977$ to $1.055$. These timings include source canonicalization, the
unchanged unweighted search and proof replay, so they do not isolate the
linear mass summation alone.

The second endpoint is complete recognition with the native stage enabled
and earlier shortcuts disabled. Both arms use the current recognizer,
switching only weighted production and verification. There is no local
native guard cap, and both retain the twelve-second global deadline.
The three stabilized circles return raw disc proofs; the five-crossing
positive uses span optimization; the trefoil continues to the exact
Khovanov fallback. Positive result bookkeeping includes the certificate hash.

\begin{center}
\small
\begin{tabular}{lrrrrr}
\toprule
Input & Old (ms) & New (ms) & Old/new & Old A/A & New A/A \\
\midrule
circle-8 & 271.428 & 212.207 & 1.279 & 1.037 & 0.992 \\
circle-16 & 855.262 & 659.084 & 1.313 & 1.001 & 1.009 \\
circle-32 & 2930.303 & 2164.070 & 1.346 & 0.999 & 1.007 \\
Optimized positive & 227.187 & 207.367 & 1.106 & 1.011 & 0.966 \\
Trefoil & 79.395 & 74.039 & 1.081 & 1.022 & 0.995 \\
\bottomrule
\end{tabular}

\end{center}

All one hundred measured recognitions and twenty warmups finish, in
101.929 seconds of driver elapsed time. Paired gains range from $1.081$
to $1.346$, with old A/A ratios $0.999$--$1.037$ and new A/A ratios
$0.966$--$1.009$. The 32-crossing circle's medians fall from 2.930303 to
2.164070 seconds, including independent positive replay. Its default
two-million-guard run remains capped; this timing does not imply a new
default-budget success. The gains concern this explicitly exercised native
workload. No ordinary-portfolio speedup or general recognition bound follows.


Reproduction uses \path{../fast/normal_orbit_research/single_orbit.py}.

Its modes are \code{audit}, \code{weighted} and \code{recognize}. The
baseline is commit \code{8160ec448}. All 475 runtime, driver, fixture and
prior-audit source pins agree across the final experiments. The incoming
archive inventory and the \code{openai/math} head \code{fd4aeeb2e} are
unchanged at this intake check. The arithmetic observation here supplements
the integrated weighted-component machinery without assuming any new
global theorem about unknot recognition.

